\section{PersonaScienceBench}
\label{sec:benchmark}

\subsection{Task and Interaction Formulation}
A \emph{task} is a scientific problem drawn from a source benchmark,
retaining that source's native deliverable and grading contract. An
\emph{actor} is the agent under evaluation; it holds the tools, executes
code where the track requires it, and produces the deliverable. A
\emph{simulated researcher} is a language model that plays the collaborating
scientist: it reviews the actor's proposed approach, answers questions,
requests revisions, and decides when the work may conclude. A \emph{persona
condition} assigns that researcher either a complete rendered profile or, in
the control condition, no identity block. An \emph{episode} is one complete
interaction from the original task statement to a sealed deliverable. An
\emph{evaluator} is the source benchmark's native grader, which runs after
the episode is sealed and never participates in the conversation.

Three information domains are kept separate by construction. \emph{Public
task context} is supplied identically in every condition and is owned by the
task, not the persona; it includes the scientific question, the input files,
and the required output format. \emph{Researcher-held information} is what the
simulated researcher can supply on request, which creates the reason to
interact at all. \emph{Private references} are the evaluator's answer keys and
grading rubrics, which are withheld from both actor and researcher throughout.
Because study-specific context is task-owned and held fixed, a persona changes
how that context is communicated, prioritized, and reviewed, but not what
ground truth exists or what information is in principle obtainable. This is
what licenses reading a score difference as a consequence of collaboration
rather than of unequal information endowment across conditions.

\figwide{1.3in}{benchmark-overview}{Placeholder: benchmark construction, context separation, and the actor--researcher--evaluator workflow.}

\subsection{Scientific Tasks and Source Conversion}
Conversion is governed by a source-preserving contract: each task keeps its
native deliverable, its native grader, and its native notion of success.
Tasks are not rewritten into a common answer format, because doing so would
discard the source-specific scoring that makes a persona effect legible as a
scientific error rather than as a conversational difference.

Four sources use an execution or design track under a shared active-researcher
protocol. DiscoveryBench tasks require the actor to produce a hypothesis
specified as context, variables, and relationship, given datasets, metadata,
and a discovery goal; its real-world split is manually derived from published
papers, while a larger synthetic split is model-generated for controlled
difficulty \citep{majumder2024discoverybench}. We retain its native
hypothesis-match grading, and note that the analysis workflow is an optional
output that the native metric does not score, so a workflow produced during
collaboration is retained as evidence rather than graded. BixBench tasks place
the actor in an empty notebook inside a provisioned environment and require
multi-step exploration of heterogeneous input files to answer expert-authored
questions; the primary regime is open answer, judged binary-correct against a
reference answer rather than by numeric tolerance
\citep{mitchener2025bixbench}.

ASI-Bench requires scientific implementation and execution scored by per-task
programmatic gates and weighted scorers against a reference solution, and it
natively varies \emph{how much methodological guidance the task supplies},
through matched levels ranging from a full method statement down to an
objective alone \citep{zhou2026asibench}. Our conversion exploits that
structure directly, and it is the source in which this benchmark's mechanism is
most literally instantiated: the actor receives the objective-only prompt,
while the withheld methodological guidance from the next-richer level is
assigned to the \emph{researcher} as optional advice the actor can obtain only
by asking. Guidance facts are extracted only from explicitly named
methodological sections and the richest guidance level is never transferred
wholesale, so the researcher holds method-level consultation rather than a
complete solution. The resulting condition is objective-only execution with
methodological consultation available on request, which is by construction not
autonomous objective-only execution, and ASI-Bench numbers reported here are
therefore not comparable to native autonomous results at that level. Because
the frozen panel draws 40 tasks from a benchmark of roughly 60, source coverage
is substantially higher here than for the other four sources.

EXP-Bench natively requires an agent to implement \emph{and execute} a research
experiment reconstructed from a published paper, with ground truth spanning
design, implementation, and conclusion components
\citep{kon2025expbench}. We convert it into a \textbf{design-only} track, and
this is a modification of the benchmark rather than a mode it provides: we
retain only its design ground truth and its design-correctness criterion, and
the actor produces an experimental design and controlled comparison plan
without executing the experiment. Execution-dependent components of the native
ground truth are therefore not scored in our conversion, and EXP-Bench results
reported here are not comparable to native EXP-Bench leaderboard numbers. Our
configuration additionally averages three judge passes per deliverable to
reduce grading variance; these are repetitions of \emph{our} grading
configuration, not part of the native EXP-Bench protocol, and they are not
independent episode trials.

SciConvBench is preserved as a dialogue track rather than being converted into
an artifact task. It retains its native ambiguous or internally inconsistent
request and its intended-specification split: the actor asks clarification
questions or returns a completion signal, with up to ten question turns and a
final attempt, and native judges assess specification recovery, intent
preservation, and conversational grounding. No scientific solver is executed.
Converting SciConvBench into a full-specification copying test would have
removed exactly the interactive difficulty it was built to measure.

\paragraph{Source pinning and leakage control.}
Every source is pinned by revision, and the datasets that ship reference
material alongside inputs require explicit leakage control. Source repositories
and dataset revisions are recorded by commit hash with per-file SHA-256
checksums, because at least one source has been revised after publication in
ways that change measured performance: the BixBench dataset was substantially
re-reviewed after its paper appeared, with questions revised for answerability
and the schema flattened, so scores are comparable only within a pinned
revision and results here are not comparable to the numbers reported in the
original paper \citep{mitchener2025bixbench}. BixBench capsule archives also
bundle the analyst's executed solution notebook alongside the input data, which
would leak answers directly into the actor's workspace. Our conversion mounts
only data members: files whose path components match reference-material
patterns are excluded, gold notebooks remain in the immutable archive outside
both actor and researcher mounts, and the researcher-visible context is
generated from input headers, labels, formats, and identifier examples without
reference to any gold notebook, hypothesis, result, or ideal answer. Where a
native grader does not verify an execution artifact, this is recorded rather
than assumed.

\subsection{Sampling and Scientific Coverage}
The task panel is a frozen random sample of 200 tasks, 40 from each source,
drawn with seed 20260929. The selection manifest is content-addressed by
SHA-256 \texttt{74619131\allowbreak 2aa95b39\allowbreak a1b04117\allowbreak
1b14e0a9\allowbreak 2dd97d7c\allowbreak b68f6d39\allowbreak d719d6bc\allowbreak
6fd20b06\allowbreak 533} and the identical panel is used for every condition and
every trial; tasks are never resampled between trials.

The sampling policy admits numerical scientific workflows and permits
diagnostic plots as intermediate artifacts, while excluding image-centric
objectives such as image interpretation, microscopy analysis, reconstruction,
computer vision, video, and visual generation. Screening was applied to
scientific objectives and required outputs, not only to input file types.
Eligibility screening and successful input preparation do not guarantee
successful execution or grading, and tasks that fail technically remain
visible as outcomes rather than being replaced.

Equal per-source quotas define the design. The resulting distribution over
scientific super-families is an empirical consequence of that design, not a
field quota: 50 life sciences, 63 physical sciences and engineering, 55
mathematics and computer science, and 32 social sciences and humanities. By
answer format the panel contains 80 empirical analyses, 40 designs, 40
specification dialogues, and 40 scientific implementations.

Source and field are therefore partly confounded, and this constrains
interpretation throughout. All selected social-science and humanities tasks
come from DiscoveryBench, and BixBench dominates the life sciences. Field
effects and format effects cannot be fully separated in this panel;
comparisons within a single source and format are reported where coverage
permits. Subfield, source-study, and task-family tags are retained in the
manifest, and multiple tasks may derive from a single source study, so tasks
are not treated as independent studies for uncertainty estimation.

\begin{table}[t]
  \centering
  \caption{Task allocation in the frozen panel. Equal per-source quotas define the design; super-family counts are the empirical distribution of the frozen sample rather than field quotas.}
  \label{tab:task-coverage}
  \begin{tabularx}{\columnwidth}{@{}lr>{\raggedright\arraybackslash}X@{}}
    \toprule
    Benchmark & Tasks & Native deliverable \\
    \midrule
    DiscoveryBench & 40 & Hypothesis (context, variables, relationship) \\
    BixBench & 40 & Open answer from notebook analysis \\
    EXP-Bench & 40 & Experimental design (design-only conversion) \\
    SciConvBench & 40 & Specification recovery dialogue \\
    ASI-Bench & 40 & Execution with researcher-held guidance \\
    \midrule
    Total & 200 & \\
    \bottomrule
  \end{tabularx}
\end{table}

\subsection{Researcher Profiles and Selection}
Profiles are complete rendered MatrAIx persona records
\citep{li2026matraix}, injected verbatim into the researcher simulator's
system identity. No truncation, summarization, or projection onto the task
domain is applied, and the same profile text is used for every task and every
trial. Neither the actor nor the evaluator receives the profile.

Selection proceeded in two stages, and the second is outcome-dependent. In the
first stage, an original sample of 300 profiles was rescreened to 180 eligible
records, from which a source-balanced, behaviorally diverse cohort of 30 was
selected \emph{without} reference to model scores. An exploratory pilot then
ran all 30 profiles plus the persona-off control. In the second stage, ten
profiles were selected by \emph{largest absolute deviation} from the
persona-off baseline in that pilot---not by highest score---giving, in ranked
order, 046, 003, 004, 089, 001, 127, 262, 224, 148, and 026. This panel
deliberately contains deviations in both directions. Persona 155 is retained as
an additional comparison condition outside the top-ten set and is labeled
separately in all analyses.

This selection is disclosed rather than presented as representative. Because
the panel was chosen using pilot outcomes, it characterizes the \emph{range} of
collaborator effects and must not be read as an unbiased estimate of an average
effect over any researcher population; the pilot scores that drove selection are
not evidence that these profiles are optimal on new tasks. Profiles are not
reselected using results from the expanded campaign.

\todo{Add the compact persona table: for each of the ten selected profiles plus
155, give domain background, methodological familiarity, communication language,
and observable workflow preferences, with links to the exact injected prompt
text and profile hashes. Two illustrative entries are available now---148 is a
senior medical researcher with medicine and statistics background, English
communication, extensive detail, accuracy-first working style, and demanding
approval behavior; 155 has an agriculture/agronomy and management background,
Mandarin communication, concrete detailed guidance, structured teamwork, and a
preference for familiar methods---but the remaining nine must be described from
their exact injected text rather than paraphrased from selection metadata.}

Two cautions apply to reading these profiles. Full profiles contain many
attributes and are sometimes internally inconsistent, so observable behavior
must be evaluated against the actual injected text rather than against a
demographic label or an assumed expertise level. And because complete profiles
are injected, communication style, availability, and domain knowledge covary
within a condition; this study does not isolate individual traits or
demographic effects (\secref{sec:limitations}).

\subsection{Collaboration Protocol and Isolation}
The four execution and design tracks share an active-researcher protocol in
which approach review and result review are mandatory. An episode proceeds in
five stages: the researcher requests input inspection and a proposed approach;
the actor reports a plan with evidence, assumptions, and uncertainties; the
researcher approves the approach or requests a focused clarification; the actor
performs the analysis or writes the design and submits a result handoff; and
the researcher requests an identifiable revision or concludes, after which the
controller seals the deliverable for separate native grading. Each episode
allows twelve researcher requests and sixty actor calls. Missing submissions
remain explicit outcomes and are not silently excluded.

Because plan and result handoffs are mandatory in this interface, handoffs must
not be counted as spontaneous engagement. Extra review requests measure
\emph{simulator} activity; questions or new proposals raised inside actor
handoffs measure \emph{actor} initiative. These are reported as separate
quantities. Shared handoffs contain actor-selected evidence only: the
researcher does not automatically inspect private tool logs, so what the
researcher knows about the analysis depends on what the actor chose to
surface.

One protocol asymmetry is recorded because it affects interpretation. The
controller permits termination only during result review, so an unreadable or
malformed input can trap a cautious researcher in approach review. This
behavior is preserved rather than patched mid-campaign, since altering it would
require a new run version and break comparability with completed episodes.

SciConvBench retains its native workflow, with the persona conditioning the
simulated user's answers. One consequence is analytically important: the
actor's \emph{first} question necessarily precedes any persona-conditioned
reply, so changes in actor activity attributable to the persona must be
measured only after the first researcher response.

Isolation is enforced on three boundaries. The persona reaches only the
researcher simulator. Private evaluator references, hidden graders, and
credentials remain outside both actor and simulator access. General knowledge
is available to the simulator in all conditions, including persona-off, so the
control differs from the persona conditions only by the presence of the
identity block. Episodes execute in isolated environments, and full shared
messages, model requests and responses, tool events, submissions, native
grades, and termination reasons are recorded for every episode, including
failed ones. Legacy optional-question bridge interfaces remain implemented for
compatibility but constitute a different interaction condition and are not used
here.
